
\documentclass{article}
\usepackage{colortbl}
\usepackage{makecell}
\usepackage{multirow}
\usepackage{supertabular}

\begin{document}

\newcounter{utterance}

\centering \large Interaction Transcript for game `ifeval', experiment `ifeval', episode 54 with Bentel rockers\_\_qwen3.5{-}4b\_\_Bentel\_iter.
\vspace{24pt}

{ \footnotesize  \setcounter{utterance}{1}
\setlength{\tabcolsep}{0pt}
\begin{supertabular}{c@{$\;$}|p{.15\linewidth}@{}p{.15\linewidth}p{.15\linewidth}p{.15\linewidth}p{.15\linewidth}p{.15\linewidth}}
    \# & \multicolumn{2}{c}{Player} && \multicolumn{2}{c}{Game Master} \\
    \hline

    \theutterance \stepcounter{utterance}  
    & & & \multicolumn{4}{p{0.6\linewidth}}{
        \cellcolor[rgb]{0.9,0.9,0.9}{
            \makecell[{{p{\linewidth}}}]{
                \texttt{\tiny{[P1$\langle$GM]}}
                \texttt{Take the text below as a starting point, and make it a complete article: "You may have to meet with a helper to work out a parenting plan. The first would be to estimate how many time you have everyday for parenting, and is that enough...."} \\
\texttt{Avoid using the following keywords: sleep, cook, feed} \\
\texttt{Mention the keyword 'schedule' for more than 5 times.} \\
            }
        }
    }
     \\ \\

    \theutterance \stepcounter{utterance}  
    & \multicolumn{4}{p{0.6\linewidth}}{
        \cellcolor[rgb]{0.561,0.737,0.561}{
            \makecell[{{p{\linewidth}}}]{
                \texttt{\tiny{[P1$\rangle$GM]}}
                \texttt{\# Navigating Parenting Responsibilities: The Importance of a Collaborative Plan} \\
\\ 
\texttt{You may have to meet with a helper to work out a parenting plan. The first step in this process is to estimate how many hours you have every day for parenting and determine if that time is sufficient for your child's needs. Without a clear understanding of your available time, it is difficult to create a balanced **schedule** that supports both your well{-}being and your child's development.} \\
\\ 
\texttt{Once you have calculated your daily hours, the next phase involves drafting a detailed **schedule** that outlines when specific tasks will occur. This document should not only list work hours but also include dedicated blocks for play, learning, and rest. A well{-}structured **schedule** helps prevent burnout by ensuring that every aspect of the child's routine is accounted for without overloading any single caregiver.} \\
\\ 
\texttt{If you are co{-}parenting or working with a support system, sharing this **schedule** becomes crucial. Both parties need to see the same timeline to avoid confusion and ensure consistency. When everyone adheres to the agreed{-}upon **schedule**, transitions become smoother, and the child feels a sense of security knowing what to expect throughout the day.} \\
\\ 
\texttt{Furthermore, regularly reviewing and adjusting the **schedule** is essential as your child grows and their needs evolve. What works today might not work tomorrow, so flexibility within the framework of a solid **schedule** is key. By maintaining a dynamic **schedule**, you can adapt to new challenges, such as school changes or extracurricular activities, while keeping the core structure intact.} \\
\\ 
\texttt{Ultimately, creating a functional parenting plan requires honest assessment and open communication. By starting with a realistic time estimate and building a comprehensive **schedule** around it, you set the foundation for a healthy and sustainable family environment. Remember, a good **schedule** is not just a list of times; it is a roadmap for success that guides every interaction and activity.} \\
            }
        }
    }
    & & \\ \\

    \theutterance \stepcounter{utterance}  
    & & & \multicolumn{2}{p{0.3\linewidth}}{
        \cellcolor[rgb]{0.9,0.9,0.9}{
            \makecell[{{p{\linewidth}}}]{
                \texttt{\tiny{[GM$|$GM]}}
                {'keywords:forbidden_words': {'forbidden_words': ['sleep', 'cook', 'feed']}, 'keywords:frequency': {'relation': 'at least', 'keyword': 'schedule', 'frequency': 6}}
            }
        }
    }
    & & \\ \\

    \theutterance \stepcounter{utterance}  
    & & & \multicolumn{2}{p{0.3\linewidth}}{
        \cellcolor[rgb]{0.9,0.9,0.9}{
            \makecell[{{p{\linewidth}}}]{
                \texttt{\tiny{[GM$|$GM]}}
                \texttt{game\_result = WIN} \\
            }
        }
    }
    & & \\ \\

\end{supertabular}
}

\end{document}
